% !TEX root = ../mainfile.tex
% !TEX encoding = UTF-8 Unicode
% !TEX spellcheck = en_US

\chapter{Discussion}
\label{ch:discussion}

This chapter draws the event-study and regression results together, relates them to the
literature, and sets out their implications and limits.

\section{The Dollar's Safe Haven Is Conditional on Shock Type}
\label{sec:disc_hierarchy}

The central finding is that the dollar's safe-haven response is conditioned by the
\emph{type} of shock, not merely by who causes it. The dollar appreciated around both
military/oil-supply shocks---the external Ukraine invasion and the U.S.-originated Hormuz
crisis---and the formal difference test cannot distinguish the two responses, even though
the United States authored one and merely witnessed the other. Yet it depreciated around the
one purely economic-policy shock it created, the Liberation Day tariffs. Origin alone does
not explain the pattern; the nature of the shock does. This refines the ``dollar-erosion''
reading \parencite{corsetti2025,kamin2025}: the safe-haven function is not weakening in
general but is specifically suspended when the United States is itself the source of an
economic-policy shock, while it survives intact---reversion notwithstanding---for the
security shocks that the safe-haven tradition has always emphasised
\parencite{ranaldo2010,habib2012}.

The regression sharpens this further. Estimated over the full sample, the dollar's
sensitivity to risk is significantly positive on normal days---the textbook safe-haven
reflex---but the interaction term shows it collapsing to essentially zero inside the tariff
window, while it remains intact during the Hormuz and Ukraine windows. The breakdown is thus
not merely a difference in the \emph{level} of the dollar across episodes but a change in its
very \emph{relationship with risk}: during the U.S.\ trade-policy shock the dollar stopped
trading like a safe haven, and only then.

\section{Gold and the Oil Channel}
\label{sec:disc_oil}

Gold does not behave as the universal safe haven that a first reading of
\textcite{baur2010} would predict. Instead it tracks the inverse of the dollar across the
U.S.-originated shocks---rising when the tariff shock weakened the dollar, falling when the
Hormuz shock strengthened it---and bids alongside the dollar only for the external Ukraine
shock. Gold is thus better described as a hedge \emph{against} the dollar than as a
substitute for it, which reinforces, rather than competes with, the chapter's main message:
the dollar's own conditional move is the organising variable, and gold prices it from the
other side.

The oil channel operates as the commodity-currency literature predicts, but only for the
supply shock. During the Hormuz crisis, oil-importer currencies underperformed
oil-exporters as crude surged, consistent with \textcite{chen2003} and \textcite{kilian2009};
during the demand-driven tariff decline the sorting was muted. The yen illustrates the
resulting tension: a classical safe haven that is also a large oil importer, it is pulled in
opposite directions during an oil-supply shock, which is part of why the basket responses
around Hormuz are smaller and noisier than the dollar-index response.

\section{Policy Implications}
\label{sec:disc_policy}

Three implications follow. For reserve managers and central banks, the dollar's safe-haven
status cannot be treated as unconditional: hedging programmes calibrated on the assumption
that the dollar always rallies in risk-off episodes will misfire precisely when the shock
originates in U.S.\ economic policy. For portfolio managers, diversification should be
thought of across \emph{shock types}, not only across assets---gold and competing-haven
currencies, not the dollar, carried the flight to quality in the tariff episode. For the
broader debate on the dollar's reserve role \parencite{obstfeld2023,jiang2026}, the evidence
is a caution against over-reading any single episode: the 2025 tariff shock did dent the
dollar, but the 2026 security shock showed the haven bid still firmly in place.

\section{Limitations}
\label{sec:disc_limitations}

Several caveats bound the results. The two most recent events---the May~2026 U.S.\ strikes
and the June~2026 ceasefire---have too few post-event trading days for the longer windows,
so they are reported only at short horizons. The regression's crisis windows are defined
around the events rather than by a structural break test, and the VIX--dollar link it
measures is a reduced-form association, not a causal channel. Gold's exceptional 2025--26
trend makes its constant-mean CARs
sensitive in level, so they are read for sign rather than magnitude. As a single-country
event study around a small number of recent, partly overlapping episodes, the design trades
breadth for the clean U.S.-versus-external contrast; and without forward/basis data beyond
the one-month tenor, covered-interest-parity deviations are left for future work.
